\documentclass[11pt]{article}

\usepackage[utf8]{inputenc}
\usepackage[T1]{fontenc}
\usepackage{geometry}
\usepackage{amsmath}
\usepackage{amssymb}
\usepackage{booktabs}
\usepackage{hyperref}
\usepackage{xurl}
\usepackage{xcolor}
\usepackage{listings}
\usepackage{enumitem}
\usepackage{parskip}
\usepackage{graphicx}
\graphicspath{{../images/}}

\geometry{margin=1in}

\lstset{
  language=Java,
  basicstyle=\ttfamily\small,
  keywordstyle=\color{blue},
  commentstyle=\color{gray},
  stringstyle=\color{teal},
  showstringspaces=false,
  columns=fullflexible,
  frame=single,
  framerule=0.2pt,
  breaklines=true
}

\setlist[itemize]{topsep=0.3em,itemsep=0.2em}
\setlist[enumerate]{topsep=0.3em,itemsep=0.2em}

% Simple per-section source line.
\newcommand{\notesources}[1]{\par\noindent{\small\color{gray}\textit{Sources:} \sloppy #1\par}}

% Unit-wise source strings for this subject (used by slides/notes).
% Path is relative to the lecture `latex/` folder (the compilation CWD).
% OOPS with Java (CSEG1044) — unit-wise sources
%
% These are short, student-facing reference strings intended to be used with:
%   - \slidesources{...}  (in beamer slides)
%   - \notesources{...}   (in lecture notes)
%
% Style inspiration: other subjects in My-Subjects/ use semicolon-separated sources
% and include an "accessed" date for web documentation.

\newcommand{\OOPSUnitISources}{Schildt, \textit{Java: A Beginner's Guide} (10e), McGraw-Hill (Java fundamentals + OOP basics).; Oracle Java Tutorials: Getting Started + Language Basics + Classes and Objects (accessed \today).; Java SE API docs (e.g., \texttt{String}, \texttt{Scanner}) (accessed \today).}

\newcommand{\OOPSUnitIISources}{Schildt, \textit{Java: The Complete Reference} (12e), McGraw-Hill (classes, inheritance, packages, interfaces).; Bloch, \textit{Effective Java} (3e), Addison-Wesley, 2018 (design choices; interfaces vs abstract classes).; Oracle Java Tutorials: Classes and Objects + Interfaces + Packages (accessed \today).; Java SE API docs (e.g., \texttt{Comparable}, \texttt{Runnable}) (accessed \today).}

\newcommand{\OOPSUnitIIISources}{Schildt, \textit{Java: The Complete Reference} (12e) (nested/inner classes, exceptions, threads, I/O).; Oracle Java Tutorials: Nested Classes + Exceptions + Concurrency + Basic I/O (accessed \today).; Java SE API docs (e.g., \texttt{Thread}, \texttt{AutoCloseable}) (accessed \today).}

\newcommand{\OOPSUnitIVSources}{Schildt, \textit{Java: The Complete Reference} (12e) (generics, lambdas, Swing, JDBC).; Bloch, \textit{Effective Java} (3e) (generics + lambdas).; Oracle Java Tutorials: Generics + Lambda Expressions + Swing + JDBC Basics (accessed \today).; Java SE API docs (e.g., \texttt{java.util.function}, \texttt{javax.swing}, \texttt{java.sql}) (accessed \today).}

\newcommand{\OOPSUnitVSources}{Schildt, \textit{Java: The Complete Reference} (12e) (Collections Framework + wrapper classes).; Oracle Java docs: Collections Framework overview (accessed \today).; Java SE API docs (\texttt{java.util} collections; wrapper classes in \texttt{java.lang}) (accessed \today).}

\newcommand{\OOPSUnitVISources}{Git SCM: \textit{Pro Git} (2e) (version control workflow), accessed \today.; GitHub Docs (repositories, issues, pull requests), accessed \today.; Oracle Java Tutorials: Swing + JDBC (accessed \today).; JUnit 5 User Guide (accessed \today).; Apache Maven Project: Getting Started (accessed \today).; Gradle User Manual: Getting Started (accessed \today).; UML class diagrams: UML 2.x overview/spec (accessed \today).}

\newcommand{\OOPSMiscWhyInterfacesSources}{Schildt, \textit{Java: The Complete Reference} (12e) (interfaces).; Bloch, \textit{Effective Java} (3e) (prefer interfaces where appropriate).; Oracle Java Tutorials: Interfaces (accessed \today).; Java SE API docs (e.g., \texttt{Comparable}, \texttt{Runnable}) (accessed \today).}



\title{Object Oriented Programming (CSEG1044)\\Unit 03 -- Lecture 04 Notes\\Java I/O Streams + FileReader/FileWriter}
\begin{document}
\maketitle
\tableofcontents

\section{Lecture Snapshot}
\begin{itemize}
  \item \textbf{Unit focus:} Nested Classes, Exceptions, Multithreading, and I/O
  \item \textbf{Main new ideas:} Stream model (big picture), Byte streams vs character streams, Text file I/O with \texttt{FileReader}/\texttt{FileWriter}
  \item \textbf{Course thread:} Use Java I/O Streams + FileReader/FileWriter to make Campus Resource Manager safer, more modular, and easier to recover when something goes wrong.
\end{itemize}
\notesources{\OOPSUnitIIISources}

\section{Prerequisites}
You should already know:
\begin{itemize}
  \item basic exception handling (\texttt{try/catch}),
  \item loops and string trimming/parsing,
  \item how to run a Java program from an IDE or terminal.
\end{itemize}

\section{Learning Outcomes}
After this lecture, you should be able to:
\begin{itemize}
  \item Explain stream model; choose byte vs character streams.
  \item Read/write text files with \texttt{FileReader}/\texttt{FileWriter} and buffering.
  \item Use \texttt{try-with-resources} for exception-safe I/O code.
\end{itemize}

\section{Suggested Reading Order}
\begin{enumerate}
  \item Read the \textbf{Prerequisite Bridge} first and confirm each recalled idea.
  \item Study the central explanation in this order: \textit{Stream model (big picture)} $\rightarrow$ \textit{Byte streams vs character streams} $\rightarrow$ \textit{Text file I/O with \texttt{FileReader}/\texttt{FileWriter}} $\rightarrow$ \textit{Reading patterns (three common styles)} $\rightarrow$ \textit{Exception-safe I/O: \texttt{try-with-resources}}.
  \item Trace the worked example line by line before checking short answers.
  \item Attempt the practice section without looking at the solution immediately.
  \item Finish with the exit questions and further-study suggestions to confirm retention.
\end{enumerate}
\notesources{\OOPSUnitIIISources}

\section{Lecture Overview}
In almost every real application you must read/write data:
\begin{itemize}
  \item read configuration,
  \item read input files (CSV/text),
  \item write reports/logs,
  \item store results for later use.
\end{itemize}
Java provides a stream-based I/O model. In this lecture we focus on text files with \texttt{FileReader}, \texttt{FileWriter}, and buffering. We also address common beginner mistakes such as wrong paths, unclosed files, and exception confusion.
\notesources{\OOPSUnitIIISources}

\section{Prerequisite Bridge}
Before reading the new material in depth, do a fast recall of the older ideas listed below.
\begin{itemize}
  \item \textbf{Exceptions}: File work is one of the most common places where exceptions matter. Main reference: \texttt{Unit 03 / Lecture 02 slides/notes}.
  \item \textbf{Strings and loops}: Reading and writing files still depends on iterating and representing text correctly. Main reference: \texttt{Unit 01 / Lecture 04 slides/notes}.
\end{itemize}
This lecture only revisits those ideas briefly so that the main explanation can stay focused on the new concept sequence.
\notesources{\OOPSUnitIIISources}

\section{Key Terms (Quick)}
\begin{itemize}
  \item \textbf{Stream}: sequence of data flowing from a source to a destination.
  \item \textbf{Input stream} vs \textbf{output stream}: read vs write.
  \item \textbf{Byte stream}: reads/writes raw bytes (\texttt{InputStream}/\texttt{OutputStream}).
  \item \textbf{Character stream}: reads/writes characters (\texttt{Reader}/\texttt{Writer}).
  \item \textbf{Buffering}: collecting data in memory to reduce expensive I/O operations.
  \item \textbf{try-with-resources}: \texttt{try(...)} that auto-closes resources.
\end{itemize}

\section{Detailed Notes}
\subsection{Stream model (big picture)}
I/O in Java is organized as streams:
\begin{itemize}
  \item A \textbf{source} (keyboard, file, network),
  \item A \textbf{stream object} (API you use),
  \item A \textbf{destination} (screen, file, network).
\end{itemize}
Streams are \textbf{one-directional}: you create an input stream to read and an output stream to write.
You can also \textbf{chain} streams: for example, a file reader wrapped by a buffered reader.

\subsection{Byte streams vs character streams}
\textbf{Byte streams} (\texttt{InputStream}/\texttt{OutputStream}) are used for binary data:
\begin{itemize}
  \item images, PDFs, audio/video, serialized binary formats.
\end{itemize}
\textbf{Character streams} (\texttt{Reader}/\texttt{Writer}) are used for text:
\begin{itemize}
  \item plain text files, logs, reading lines of text.
\end{itemize}
\paragraph{Rule of thumb.} Use \texttt{Reader/Writer} for text; use \texttt{InputStream/OutputStream} for binary.

\subsection{Text file I/O with \texttt{FileReader}/\texttt{FileWriter}}
Basic idea:
\begin{itemize}
  \item \texttt{FileReader} reads characters from a file.
  \item \texttt{FileWriter} writes characters to a file.
  \item These can be wrapped in buffered versions for performance.
\end{itemize}

\paragraph{Important note (encoding).}
\texttt{FileReader/FileWriter} use the platform default charset.
For most course examples this is fine, but in real projects you often choose a charset explicitly (UTF-8).

\textbf{Buffering:}
\begin{itemize}
  \item \texttt{BufferedReader} provides \texttt{readLine()} and reads in chunks.
  \item \texttt{BufferedWriter} provides efficient writing and \texttt{newLine()}.
\end{itemize}

\subsection{Reading patterns (three common styles)}
\textbf{1) Read character-by-character} (rarely used in app code):
\begin{itemize}
  \item simple idea, but slower and harder to manage.
\end{itemize}

\textbf{2) Read line-by-line} (most common for text files):
\begin{itemize}
  \item use \texttt{BufferedReader.readLine()}.
\end{itemize}

\textbf{3) Read tokens/CSV} (usually needs parsing/splitting after reading lines):
\begin{itemize}
  \item read line $\rightarrow$ split by delimiter $\rightarrow$ parse fields.
\end{itemize}

\subsection{Exception-safe I/O: \texttt{try-with-resources}}
File I/O can fail for many reasons: wrong path, missing file, permissions.
Also, resources must be closed.

Use:
\begin{itemize}
  \item \texttt{try (BufferedReader br = ...)\{...\}} to auto-close even if exception occurs.
\end{itemize}
This reduces resource leaks and makes code cleaner.

\subsection{Paths and a common beginner mistake}
\textbf{The JVM reads files relative to the ``working directory''} (where you run \texttt{java}).
If the file is ``not found'':
\begin{itemize}
  \item confirm the path,
  \item print the working directory (or use absolute path),
  \item ensure file extension and capitalization are correct.
\end{itemize}
\paragraph{Quick debug snippet (print working directory).}
\begin{lstlisting}
System.out.println(new java.io.File(".").getAbsolutePath());
\end{lstlisting}

\paragraph{Windows path gotcha (Java strings).}
In Java string literals, backslash is an escape character:
\begin{itemize}
  \item wrong: \verb|"D:\data\marks.txt"| (contains escape sequences)
  \item correct: \verb|"D:\\data\\marks.txt"| or \verb|"D:/data/marks.txt"|
\end{itemize}
\subsection{Overwrite vs append (writing)}
\texttt{FileWriter(path)} overwrites the file by default.
To append:
\begin{lstlisting}
new FileWriter("log.txt", true); // append=true
\end{lstlisting}
\notesources{\OOPSUnitIIISources}

\section{Running Project Connection}
Use Java I/O Streams + FileReader/FileWriter to make Campus Resource Manager safer, more modular, and easier to recover when something goes wrong.
\begin{itemize}
  \item Use Campus Resource Manager as the running case study, or map the same idea to your own project domain if you are using another example.
  \item Ask where today's idea should live: model, service, DAO, UI, or team workflow.
  \item Use the lecture example as a smaller version of the same design choice you will make in the capstone.
\end{itemize}
\notesources{\OOPSUnitIIISources}

\section{Common Mistakes and Confusions}
\begin{itemize}
  \item Forgetting to close streams or readers.
  \item Using byte streams when character streams fit text better.
  \item Reporting vague file errors with no recovery hint.
\end{itemize}
\notesources{\OOPSUnitIIISources}

\section{Worked Example (tutorial): read marks and write a report}
\subsection*{Goal}
Read integer marks from \texttt{marks.txt} (one mark per line), compute sum/avg/max, and write a report to \texttt{report.txt}.

\subsection*{Input file example}
\begin{lstlisting}[language=]
78
91
65

88
\end{lstlisting}
Blank lines are allowed; we will ignore them. If the file has no numeric rows, report the maximum as \texttt{N/A} instead of printing a fake value.

\subsection*{Code}
\begin{lstlisting}
import java.io.BufferedReader;
import java.io.BufferedWriter;
import java.io.FileReader;
import java.io.FileWriter;
import java.io.IOException;

public class MarksReport {
    public static void main(String[] args) {
        String inPath = "marks.txt";
        String outPath = "report.txt";

        int sum = 0, max = 0, count = 0;

        try (BufferedReader br = new BufferedReader(new FileReader(inPath))) {
            String line;
            while ((line = br.readLine()) != null) {
                line = line.trim();
                if (line.isEmpty()) continue;

                int x = Integer.parseInt(line);
                if (count == 0 || x > max) max = x;
                sum += x;
                count++;
            }
        } catch (IOException e) {
            System.out.println("I/O error: " + e.getMessage());
            return;
        } catch (NumberFormatException e) {
            System.out.println("Invalid number in input file: " + e.getMessage());
            return;
        }

        double avg = (count == 0) ? 0.0 : (sum / (double) count);

        try (BufferedWriter bw = new BufferedWriter(new FileWriter(outPath))) {
            bw.write("Count: " + count);
            bw.newLine();
            bw.write("Sum: " + sum);
            bw.newLine();
            bw.write("Max: " + (count == 0 ? "N/A" : max));
            bw.newLine();
            bw.write(String.format("Avg: %.2f", avg));
            bw.newLine();
        } catch (IOException e) {
            System.out.println("Write error: " + e.getMessage());
        }
    }
}
\end{lstlisting}

\subsection*{What to notice}
\begin{itemize}
  \item \texttt{try-with-resources} closes readers/writers automatically.
  \item \texttt{BufferedReader.readLine()} makes line-by-line reading easy.
  \item We handle both I/O errors and invalid numbers separately (better messages).
  \item We use \texttt{String.format} for formatted average.
\end{itemize}

\subsection*{Common improvements (good practice)}
\begin{itemize}
  \item Validate mark range (0--100) and count invalid rows.
  \item Show a clearer message when the file is missing (path help).
  \item Write results in a CSV format if needed.
\end{itemize}

\section{Demo / Observation Task}
Use this block as a short reproducible demo or observation task.
\begin{itemize}
  \item Reproduce the lecture's main example around \textit{Stream model (big picture)} once without looking at the solution first.
  \item Change one input, rule, or design choice in Campus Resource Manager and predict the result before checking it.
  \item Record one success case, one edge case, and one failure case so the idea becomes observable instead of purely theoretical.
\end{itemize}
\notesources{\OOPSUnitIIISources}

\section{Mini tutorial: copy a text file}
Copying is a great exercise because it uses both reading and writing.
\begin{lstlisting}
import java.io.*;

public class CopyTextFile {
    public static void main(String[] args) throws IOException {
        String src = "input.txt";
        String dst = "output.txt";

        try (BufferedReader br = new BufferedReader(new FileReader(src));
             BufferedWriter bw = new BufferedWriter(new FileWriter(dst))) {

            String line;
            while ((line = br.readLine()) != null) {
                bw.write(line);
                bw.newLine();
            }
        }
    }
}
\end{lstlisting}

\section{Practice Check (with brief solutions)}
\begin{enumerate}
  \item Which should be used for text files: byte stream or character stream?\par
  \textbf{Answer:} character stream (\texttt{Reader/Writer}).

  \item Why do we use buffering?\par
  \textbf{Answer:} it reduces the number of expensive disk I/O calls by reading/writing in chunks.

  \item If a file is not found, name one common cause.\par
  \textbf{Answer:} wrong working directory / wrong path.
\end{enumerate}

\section{Self-Study Practice (with hints)}
\begin{enumerate}
  \item Write a program that copies a text file (read from one, write to another) using buffered streams.
  \textbf{Hint:} the copy tutorial above is the solution pattern.

  \item Write a program that counts lines and words in a file.
  \textbf{Hint:} count lines with \texttt{readLine()} and count words with \path|split("\\s+")|.

  \item Modify the worked example to ignore marks outside 0--100 and count how many were invalid.
  \textbf{Hint:} after parsing, check range and continue.
\end{enumerate}

\section{Summary (what to remember)}
\begin{itemize}
  \item Streams model: input vs output; chain streams for features like buffering.
  \item Use \texttt{Reader/Writer} for text and buffer with \texttt{BufferedReader/BufferedWriter}.
  \item Use \texttt{try-with-resources} for reliable closing.
  \item Most file bugs come from wrong paths/working directory: print it and verify.
  \item Handle \texttt{IOException} separately from parse/validation exceptions.
\end{itemize}

\section{Lab Alignment (recommended)}
\begin{itemize}
  \item Exp-13: File Handling using I/O streams (based on this lecture)
\end{itemize}

\section{Further Study}
\begin{itemize}
  \item Revisit the prerequisite references if any recall bullet still feels weak.
  \item Rework the lecture example with one changed assumption, input, or design choice.
  \item Read the textbook and documentation references listed in the source line for a second explanation of the same topic.
  \item Apply the same idea once inside Campus Resource Manager so that the concept moves from theory to design practice.
\end{itemize}
\notesources{\OOPSUnitIIISources}

\section{Exit Questions (with sample answers)}
\textbf{Q1.} Why is \texttt{try-with-resources} preferred?\par
\textbf{Sample answer:} It automatically closes resources even if an exception occurs, preventing resource leaks and reducing boilerplate code.

\vspace{0.6em}
\textbf{Q2.} When do you choose byte streams instead of character streams?\par
\textbf{Sample answer:} Use byte streams for binary data like images or PDFs. Use character streams for text data.

\vspace{0.6em}
\textbf{Q3.} Why is try-with-resources preferred for file work?\par
\textbf{Sample answer:} It closes resources automatically and keeps the success path easier to read.

\end{document}
